% GENERATED FILE. Edit canonical lessons, then run npm run generate:books.
% canonical-source-hash: fnv1a64:5b6c862191b9cdd5
% canonical-lessons: KA-C285-mukkalu, KA-C285-marudina, KA-C285-balya, KA-C285-amavasye, KA-C285-madhyantara

\chapter{Three quarters, Next day, Childhood, New moon day, Interval}
\label{ch:ka-three-quarters-next-day-childhood-new-moon-day-interval}
\hlchaptermodality{285}

\begin{chapteropening}
\textbf{By the end of this chapter:} \emph{I can say three quarters, the next day, childhood, the new moon day and an interval.}

Complete the last lesson of chapter 285: I can say three quarters, the next day, childhood, the new moon day and an interval.
\end{chapteropening}

\section[\texorpdfstring{\kn{ಮುಕ್ಕಾಲು} (mukkālu)}{ಮುಕ್ಕಾಲು (mukkālu)}]{\kn{ಮುಕ್ಕಾಲು} (mukkālu) — three quarters}

\label{lesson:KA-C285-mukkalu}



\begin{quote}
Before the new one: say the Kannada for time, then the Kannada for sometimes.
\end{quote}

\subsection*{You'll want to know: \kn{ಮುಕ್ಕಾಲು}}
\textbf{\kn{ಮುಕ್ಕಾಲು}} — \emph{mukkālu} — \textquotedblleft{}three quarters\textquotedblright{}.

\begin{grammarlens}[title={What you've built}]
One more word for this chapter.
\end{grammarlens}

\subsection*{Your turn}
\begin{itemize}
\raggedright
  \item \emph{Say it:} \emph{mukkālu}
  \item \emph{Say it:} \emph{mukkālu}, once more
  \item \emph{Say it:} say \emph{mukkālu}
  \item \emph{Recall:} say the Kannada for time, then the Kannada for sometimes, then say \emph{mukkālu} again
\end{itemize}

\begin{tcolorbox}[breakable,colback=teal!4,colframe=teal!35!black,title={Before you move on}]
What does \kn{ಮುಕ್ಕಾಲು} mean? (\textquotedblleft{}Three quarters\textquotedblright{}.) Say it once more.
\end{tcolorbox}
\section[\texorpdfstring{\kn{ಮರುದಿನ} (marudina)}{ಮರುದಿನ (marudina)}]{\kn{ಮರುದಿನ} (marudina) — the next day, the following day}

\label{lesson:KA-C285-marudina}



\begin{quote}
Before the new one: say the Kannada for sometimes, then the Kannada for three quarters.
\end{quote}

\subsection*{You'll want to know: \kn{ಮರುದಿನ}}
\textbf{\kn{ಮರುದಿನ}} — \emph{marudina} — \textquotedblleft{}the next day, the following day\textquotedblright{}.

\begin{grammarlens}[title={What you've built}]
One more word for this chapter.
\end{grammarlens}

\subsection*{Your turn}
\begin{itemize}
\raggedright
  \item \emph{Say it:} \emph{marudina}
  \item \emph{Say it:} \emph{marudina}, once more
  \item \emph{Say it:} say \emph{marudina}
  \item \emph{Recall:} say the Kannada for sometimes, then the Kannada for three quarters, then say \emph{marudina} again
\end{itemize}

\begin{tcolorbox}[breakable,colback=teal!4,colframe=teal!35!black,title={Before you move on}]
What does \kn{ಮರುದಿನ} mean? (\textquotedblleft{}The next day, the following day\textquotedblright{}.) Say it once more.
\end{tcolorbox}
\section[\texorpdfstring{\kn{ಬಾಲ್ಯ} (bālya)}{ಬಾಲ್ಯ (bālya)}]{\kn{ಬಾಲ್ಯ} (bālya) — childhood}

\label{lesson:KA-C285-balya}



\begin{quote}
Before the new one: say the Kannada for three quarters, then the Kannada for the next day.
\end{quote}

\subsection*{You'll want to know: \kn{ಬಾಲ್ಯ}}
\textbf{\kn{ಬಾಲ್ಯ}} — \emph{bālya} — \textquotedblleft{}childhood\textquotedblright{}.

\begin{grammarlens}[title={What you've built}]
One more word for this chapter.
\end{grammarlens}

\subsection*{Your turn}
\begin{itemize}
\raggedright
  \item \emph{Say it:} \emph{bālya}
  \item \emph{Say it:} \emph{bālya}, once more
  \item \emph{Say it:} say \emph{bālya}
  \item \emph{Recall:} say the Kannada for three quarters, then the Kannada for the next day, then say \emph{bālya} again
\end{itemize}

\begin{tcolorbox}[breakable,colback=teal!4,colframe=teal!35!black,title={Before you move on}]
What does \kn{ಬಾಲ್ಯ} mean? (\textquotedblleft{}Childhood\textquotedblright{}.) Say it once more.
\end{tcolorbox}
\section[\texorpdfstring{\kn{ಅಮಾವಾಸ್ಯೆ} (amāvāsye)}{ಅಮಾವಾಸ್ಯೆ (amāvāsye)}]{\kn{ಅಮಾವಾಸ್ಯೆ} (amāvāsye) — the new moon day}

\label{lesson:KA-C285-amavasye}



\begin{quote}
Before the new one: say the Kannada for the next day, then the Kannada for childhood.
\end{quote}

\subsection*{You'll want to know: \kn{ಅಮಾವಾಸ್ಯೆ}}
\textbf{\kn{ಅಮಾವಾಸ್ಯೆ}} — \emph{amāvāsye} — \textquotedblleft{}the new moon day\textquotedblright{}.

\begin{grammarlens}[title={What you've built}]
One more word for this chapter.
\end{grammarlens}

\subsection*{Your turn}
\begin{itemize}
\raggedright
  \item \emph{Say it:} \emph{amāvāsye}
  \item \emph{Say it:} \emph{amāvāsye}, once more
  \item \emph{Say it:} say \emph{amāvāsye}
  \item \emph{Recall:} say the Kannada for the next day, then the Kannada for childhood, then say \emph{amāvāsye} again
\end{itemize}

\begin{tcolorbox}[breakable,colback=teal!4,colframe=teal!35!black,title={Before you move on}]
What does \kn{ಅಮಾವಾಸ್ಯೆ} mean? (\textquotedblleft{}The new moon day\textquotedblright{}.) Say it once more.
\end{tcolorbox}
\section[\texorpdfstring{\kn{ಮಧ್ಯಂತರ} (madhyantara)}{ಮಧ್ಯಂತರ (madhyantara)}]{\kn{ಮಧ್ಯಂತರ} (madhyantara) — an interval (in a film, a match)}

\label{lesson:KA-C285-madhyantara}



\begin{quote}
Before the new one: say the Kannada for childhood, then the Kannada for the new moon day.
\end{quote}

\subsection*{You'll want to know: \kn{ಮಧ್ಯಂತರ}}
\textbf{\kn{ಮಧ್ಯಂತರ}} — \emph{madhyantara} — \textquotedblleft{}an interval (in a film, a match)\textquotedblright{}.

\begin{grammarlens}[title={What you've built}]
One more word for this chapter.
\end{grammarlens}

\subsection*{Your turn}
\begin{itemize}
\raggedright
  \item \emph{Say it:} \emph{madhyantara}
  \item \emph{Say it:} \emph{madhyantara}, once more
  \item \emph{Say it:} say \emph{madhyantara}
  \item \emph{Recall:} say the Kannada for childhood, then the Kannada for the new moon day, then say \emph{madhyantara} again
\end{itemize}

\begin{tcolorbox}[breakable,colback=teal!4,colframe=teal!35!black,title={Before you move on}]
What does \kn{ಮಧ್ಯಂತರ} mean? (\textquotedblleft{}An interval (in a film, a match)\textquotedblright{}.) Say it once more.
\end{tcolorbox}
